% 数据库系统原理课程设计 —— 系统设计与开发报告
% 构建：在本目录执行 make（需要 XeLaTeX、latexmk 与 Graphviz）
\documentclass[UTF8,zihao=-4,fontset=mac,a4paper]{ctexart}

% ---------- 封面信息（提交前填写） ----------
\newcommand{\studentname}{（姓名）}
\newcommand{\studentid}{（学号）}
\newcommand{\advisor}{（指导教师）}
\newcommand{\reportdate}{2026 年 11 月}
% --------------------------------------------

\usepackage[margin=2.5cm]{geometry}
\usepackage{graphicx,booktabs,longtable,tabularx,array,multirow,amsmath,amssymb,xcolor,caption,enumitem,float}
\usepackage{listings}
\usepackage{pgfplots}
\pgfplotsset{compat=1.18}
\usepackage[hidelinks,bookmarksnumbered]{hyperref}

\graphicspath{{build/figures/}}
\setlength{\parskip}{0.2em}
\linespread{1.3}
\setlist{nosep,leftmargin=2em}
\captionsetup{font=small,labelsep=quad}
\ctexset{
  section={format=\Large\heiti,name={第,章},number=\chinese{section}},
  subsection={format=\large\heiti},
  subsubsection={format=\normalsize\heiti},
}
\renewcommand{\thesubsection}{\arabic{section}.\arabic{subsection}}
\renewcommand{\thesubsubsection}{\thesubsection.\arabic{subsubsection}}
\renewcommand{\thefigure}{\arabic{section}-\arabic{figure}}
\renewcommand{\thetable}{\arabic{section}-\arabic{table}}
\counterwithin*{figure}{section}
\counterwithin*{table}{section}

\lstset{basicstyle=\ttfamily\footnotesize,breaklines=true,frame=tb,framesep=4pt,columns=fullflexible,
  keywordstyle=\color{blue!60!black}\bfseries,commentstyle=\color{gray},showstringspaces=false,
  xleftmargin=1em,xrightmargin=1em}
\lstdefinelanguage{PLpgSQL}[]{SQL}{morekeywords={RETURNS,LANGUAGE,TRIGGER,PERFORM,RAISE,EXCEPTION,USING,ERRCODE,
  REPLACE,FUNCTION,IF,THEN,END,RETURN,EXISTS,BEFORE,EACH,ROW,EXECUTE,GRANT,REVOKE,ROLE,OLD,NEW,WHEN,LATERAL}}

\newcommand{\code}[1]{\texttt{\small #1}}
% Long identifiers such as trg_cycle_upstream_costs_closed may break after an underscore.
\DeclareRobustCommand{\_}{\textunderscore\allowbreak}
\setlength{\emergencystretch}{2em}

\begin{document}

\begin{titlepage}
  \centering
  \vspace*{2.5cm}
  {\Large 同济大学计算机科学与技术学院\par}
  \vspace{0.8cm}
  {\Large 数据库系统原理课程设计\par}
  \vspace{2.2cm}
  {\heiti\fontsize{26pt}{36pt}\selectfont 大模型 API 用量统计\\[0.3em] 与费用分摊系统\par}
  \vspace{0.8cm}
  {\Large 系统设计与开发报告\par}
  \vspace{3cm}
  \begin{tabular}{rl}
    姓\quad 名： & \studentname \\[0.6em]
    学\quad 号： & \studentid \\[0.6em]
    指导教师： & \advisor \\[0.6em]
    数据库系统： & PostgreSQL 17（兼容 SQLite 3） \\[0.6em]
    开发语言： & Python 3.12 / JavaScript（Vue 3） \\
  \end{tabular}
  \vfill
  {\large \reportdate\par}
\end{titlepage}

\pagenumbering{Roman}
\begin{abstract}
本系统面向多人共享大模型 API 服务的场景：一组成员共同购买若干上游账号（包月订阅或按量计费的 API 渠道），通过统一的 API 网关分发个人密钥，月末需要按各人实际用量公平分摊总成本。系统管理用户、API 密钥及其历史归属、调用记录、版本化的模型价格、上游成本与梯度分摊规则，实现调用级计价、按用户/模型/时间的用量统计、共享成本分摊和月度账单生成。设计上以“账单可复现、可核对”为核心目标：价格按版本保存，计价结果记录所用版本，账期关账时固化成本快照与梯度规则，并由数据库触发器保证已关账数据不可篡改。数据库采用 PostgreSQL，逻辑模式共 32 张基本表、5 个视图，满足 3NF（快照字段为有意保留的历史事实），并通过 CHECK 约束、部分唯一索引、16 个触发器和四级数据库角色落实完整性与安全性设计。系统在 100 万条调用记录规模下完成了索引与 SQL 改写的性能测试，常用查询在 5 ms 内返回。

\noindent\textbf{关键词：}大模型 API；用量统计；成本分摊；PostgreSQL；触发器；数据库设计
\end{abstract}
\newpage

\tableofcontents
\newpage
\pagenumbering{arabic}

\section{问题定义与可行性分析}

\subsection{项目背景}

大模型 API 的计费日趋复杂：厂商既提供附带调用额度的包月订阅，也提供按 Token 计费的 API 渠道。按量计费又区分输入、输出、缓存读取、缓存写入四种单价，并随服务档位（Priority、Flex）和上下文长度变化。个人单独购买多个订阅成本高、利用率低，因此常见做法是多人合买若干上游账号，通过 API 网关（本项目对接开源网关 CLIProxyAPI，简称 CPA）向每人分发独立密钥，共享使用。

共享带来了核算问题。月度总成本由固定订阅费和按量费用组成，成员用量相差可达数十倍，订阅周期与自然月不对齐，价格经常调整，密钥还会被重置、吊销或转手。仅靠网关统计和电子表格手工结算，会出现以下问题：
\begin{enumerate}
  \item 调用记录分散在网关的 SQLite 日志库中，缺少用户维度，无法回答“某人本月在哪些模型上花了多少钱”；
  \item 价格变化后无法说明某一期账单是按哪一版价格算出来的，账单不可复核；
  \item 密钥转手后，历史调用被错误地计入新主人名下；
  \item 分摊规则（例如用量越大单价越低的累进折扣）改动后，已经结清的月份会被“追溯”改变。
\end{enumerate}

\subsection{问题定义}

本课题要建立一个以关系数据库为核心的“大模型 API 用量统计与费用分摊系统”，其目标是：
\begin{itemize}
  \item \textbf{统一管理数据}：用户、API 密钥及其历史归属、调用记录、版本化的模型价格、上游账号成本、分摊规则、账期与账单；
  \item \textbf{调用级计价}：对每一次调用按当时生效的价格版本计算等效费用，并保存计算明细；
  \item \textbf{多维统计}：支持按用户、模型、时间、上游账号等维度统计请求数、Token 与费用；
  \item \textbf{共享成本分摊}：把每个账期的上游总成本按梯度规则公平地分到每位成员，精确到分且总额守恒；
  \item \textbf{账单可查询、可核对}：每期账单记录所用价格版本、规则快照与成本快照，关账后数据由数据库保证不可修改。
\end{itemize}

系统边界：网关 CPA 负责转发调用与管理密钥，用量库 CPAMP（网关的用量记录与价格同步组件）负责记录原始调用日志和抓取公开价格。本系统以只读方式增量读取用量库，通过网关管理接口增删密钥和查询上游额度，自身不转发任何模型调用。

\subsection{可行性分析}

\paragraph{技术可行性}
按编程助手单人每天数百次调用估算，30 人左右的共享组每月约 10 万至 50 万次调用，三年累计不超过数千万行，单机 PostgreSQL 足以承载；第 7 章的测试表明百万级数据上常用查询在 5 ms 内完成。计价与分摊均为整数运算，没有浮点误差。所用框架均为成熟的开源项目。

\paragraph{经济可行性}
系统全部基于开源软件，可与网关部署在同一台 2 核 4 GB 云服务器上，不增加硬件费用，而自动分摊可以省去每月数小时的手工对账。

\paragraph{操作可行性}
成员通过已在使用的 Telegram 机器人注册和管理密钥，Web 端用自己的密钥登录；管理员的日常操作只有“创建账期—预览—关账”三步，其余由后台任务完成。

\paragraph{法律与安全可行性}
系统不存储提示词与回复内容；密钥只保存不可逆摘要和脱敏值，上游账号凭据始终由网关保管。

\subsection{开发环境与工具}

\begin{table}[H]
\centering\small
\caption{开发与运行环境}
\begin{tabularx}{\textwidth}{@{}lX@{}}
\toprule
类别 & 选型 \\ \midrule
数据库管理系统 & PostgreSQL 17（生产可选），SQLite 3.40+（单机默认，兼容） \\
后端 & Python 3.12，FastAPI，SQLAlchemy 2.1（ORM 与 Core），Alembic 1.20（模式迁移），psycopg 3 \\
前端 & Vue 3.5，Vuetify 3，ECharts 6，Vite 6，Vitest \\
交互入口 & Web 控制台（B/S），Telegram Bot，命令行工具 \\
设计工具 & Graphviz（E-R 图、数据流图），本报告使用 XeLaTeX 排版 \\
部署 & Docker Compose，GitHub Actions 持续集成（SQLite 与 PostgreSQL 双库测试） \\
\bottomrule
\end{tabularx}
\end{table}

\section{需求分析}

\subsection{用户角色}

系统的使用者按权限从低到高分为五类，另有一个无人值守的后台任务（表~\ref{tab:roles}）。

\begin{table}[H]
\centering\small
\caption{用户角色与权限}\label{tab:roles}
\begin{tabularx}{\textwidth}{@{}lXX@{}}
\toprule
角色 & 身份来源 & 主要权限 \\ \midrule
匿名访问者 & 无 & 访问登录页、健康检查 \\
只读访客 & 未绑定用户的 API 密钥登录 & 仅查看该密钥自己的请求明细与站点状态 \\
成员 & 已注册用户名下的有效密钥登录；Telegram 私聊 & 查看账务总览、排行、价格规则、个人用量；新增、命名、重置、吊销自己的密钥 \\
管理员 & 被授予管理权限的成员，或持有管理令牌者 & 全部成员权限；账期创建/预览/关账、账单调整、手工用量、改价与价格同步、梯度规则与上游成本配置、密钥转移与计费倍率、授予/撤销管理员 \\
后台任务 & 部署配置 & 增量同步调用记录、计价、同步网关密钥 \\
\bottomrule
\end{tabularx}
\end{table}

\subsection{功能需求}

\begin{longtable}{@{}p{0.17\textwidth}p{0.77\textwidth}@{}}
\caption{功能需求}\\
\toprule 编号 & 需求 \\ \midrule\endfirsthead
\toprule 编号 & 需求 \\ \midrule\endhead
F1 用户与认证 & 群成员向机器人发送 \code{/register} 注册并获得第一把密钥；Web 端以密钥或独立的管理令牌登录，登录失败限流 \\
F2 密钥管理 & 成员新增、命名、重置、吊销自己的密钥，明文只显示一次；管理员转移密钥（指定生效时间）、为未绑定的按量密钥设置倍率；每次归属变化形成带起止时间的历史 \\
F3 调用采集 & 增量读取用量库，按源记录号与事件摘要去重；无法解析的记录进入死信等待重试 \\
F4 价格与计价 & 同步、改价、重新发布都生成新价格版本；按“版本 + 模型 + 档位 + 上下文段”计价并保存明细；价格变动只重算未关账账期内变价的模型 \\
F5 成本与规则 & 上游账号分组，每组一条梯度规则；包月账号配置订阅周期与费用（按天折算），按量渠道配置汇率 \\
F6 账期与账单 & 创建账期时固化价格版本、规则与成本快照；开放期间随时预览，关账后定稿不可改；关账前可录入账单调整与手工用量 \\
F7 统计查询 & 账务总览、排行、单用户按模型与档位用量、多条件请求明细、站点状态、上游额度，支持任意时间范围 \\
F8 系统管理 & 审计日志（含前后值）、三方对账、备份、跨数据库迁移 \\
\bottomrule
\end{longtable}

\subsection{非功能需求}

\begin{itemize}
  \item \textbf{正确性}：金额全部以整数存储（分、纳美元、百万分之一），所有舍入采用四舍五入并在分摊时保证总额守恒；
  \item \textbf{可追溯}：任何一张已关账账单都能由数据库中的快照完整重算；
  \item \textbf{性能}：常用页面查询在百万级调用记录下响应时间不超过 100 ms，同步与计价每 10 秒一轮；
  \item \textbf{安全}：密钥不落明文，会话令牌只存摘要，所有写操作校验 CSRF，按角色控制数据可见范围；
  \item \textbf{可维护}：数据库结构由迁移脚本管理，支持在线备份，支持从 SQLite 平滑迁移到 PostgreSQL。
\end{itemize}

\subsection{数据流图}

图~\ref{fig:dfd0} 是系统的顶层（0 层）数据流图，系统与四个外部实体交换数据。图~\ref{fig:dfd1} 把系统分解为六个加工和六个数据存储。

\begin{figure}[H]
  \centering
  \includegraphics[width=0.92\textwidth]{dfd0.pdf}
  \caption{顶层数据流图}\label{fig:dfd0}
\end{figure}

\begin{figure}[H]
  \centering
  \includegraphics[height=0.78\textheight]{dfd1.pdf}
  \caption{1 层数据流图}\label{fig:dfd1}
\end{figure}

\subsection{数据字典}

数据字典按结构化分析的五类条目给出。各数据存储对应的完整表结构（字段、类型、约束）见第 4 章与附录 A，由程序从数据库模型自动生成，保证与实现一致。

\subsubsection{数据项（节选）}

\begin{longtable}{@{}p{0.2\textwidth}p{0.2\textwidth}p{0.18\textwidth}p{0.34\textwidth}@{}}
\caption{主要数据项}\\
\toprule 数据项 & 类型与长度 & 取值范围 & 含义与约定 \\ \midrule\endfirsthead
\toprule 数据项 & 类型与长度 & 取值范围 & 含义与约定 \\ \midrule\endhead
时间戳（*\_ms） & 64 位整数 & $\ge 0$ & Unix 毫秒时间；区间统一左闭右开 \\
等效费用（*\_nano\_usd） & 64 位整数 & $\ge 0$ & 单位 $10^{-9}$ 美元，计价结果不做中间舍入 \\
单价（*\_nano\_per\_token） & 64 位整数 & $\ge 0$ & 纳美元/Token，等于“美元/百万 Token × 1000” \\
金额（*\_cents） & 64 位整数 & 调整可为负 & 人民币分 \\
倍率与汇率（*\_ppm） & 64 位整数 & $\ge 0$ & 百万分之一定点数，如 0.9 存为 900000 \\
Token 计数 & 64 位整数 & $\ge 0$ & 输入、输出、推理、缓存读、缓存写、合计 \\
密钥摘要 & 定长字符 64 & 十六进制 & SHA-256，用于与调用记录匹配 \\
密钥状态 & 字符 20 & active、unowned、revoked、retired & 使用中/未绑定/已吊销/网关中已移除 \\
账期状态 & 字符 20 & open、preview、closed & 开放/已预览/已关账 \\
价格版本状态 & 字符 20 & draft、active、retired & 同一时刻至多一个 active \\
服务档位 & 字符 40 & default、priority、flex & fast 视同 priority \\
梯度规则 & JSON 文本 & 连续分段 & 形如 \code{[\{left:0,right:300,multiplier:1\},…]}，单位美元 \\
\bottomrule
\end{longtable}

\subsubsection{数据结构}

\begin{itemize}
  \item \textbf{调用记录} = 记录号 + 发生时间 + 模型名 + 服务档位 + 密钥摘要 + 上游账号 + \{Token 计数\}$_6$ + 是否失败 + 延迟 + 首 Token 时间；
  \item \textbf{模型价格} = 价格版本 + 模型名 + 默认四价 + [Priority 四价] + [Flex 四价] + [上下文分段价表]；
  \item \textbf{账单} = 账期 + 用户 + 实际用量 + 加权用量 + 应付金额 + 调整金额 + \{账单明细\}；
  \item \textbf{账单明细} = 资源池 + 分组 + 实际用量 + 加权用量 + 金额 + 密钥数；
  \item \textbf{账期快照} = 价格版本 + 梯度规则副本 + \{上游账号成本\} + \{分组结算结果\}。
\end{itemize}

\subsubsection{数据流}

\begin{longtable}{@{}p{0.17\textwidth}p{0.16\textwidth}p{0.16\textwidth}p{0.43\textwidth}@{}}
\caption{主要数据流}\\
\toprule 名称 & 来源 & 去向 & 组成与说明 \\ \midrule\endfirsthead
\toprule 名称 & 来源 & 去向 & 组成与说明 \\ \midrule\endhead
增量调用记录 & 用量库 & P1 & 源记录号大于检查点的调用记录，每批 1000 条 \\
未计价记录 & D1 & P2 & 生效价格版本下尚无计价结果的调用记录，每批 500 条 \\
生效价格 & D3 & P2 & 状态为 active 的价格版本及其模型价格 \\
调用时刻归属 & D4 & P2 & 满足“生效时间 $\le t <$ 失效时间”的归属期 \\
密钥操作请求 & 成员 & P3 & 操作类型 + 目标密钥 + 当前密钥复核 \\
账期操作 & 管理员 & P5 & 账期名 + 起止时间 + 规则；或预览/关账/调整指令 \\
账期用量 & D2 & P5 & 按用户 × 上游账号 × 资源池汇总的等效费用 \\
账单与明细 & P5 & D5 & 每位用户一张账单及其按分组、资源池的明细 \\
查询条件 & 成员 & P6 & 时间范围 + 用户/密钥/模型/档位/状态等筛选条件 \\
\bottomrule
\end{longtable}

\subsubsection{数据存储}

\begin{longtable}{@{}p{0.08\textwidth}p{0.2\textwidth}p{0.3\textwidth}p{0.34\textwidth}@{}}
\caption{数据存储}\\
\toprule 编号 & 名称 & 对应基本表 & 存取方式 \\ \midrule\endfirsthead
\toprule 编号 & 名称 & 对应基本表 & 存取方式 \\ \midrule\endhead
D1 & 原始调用记录 & raw\_usage\_events, dead\_letters, sync\_checkpoints & 只追加；按时间、密钥摘要检索 \\
D2 & 计价记录 & rated\_events & 只追加；按版本+时间、用户+时间聚合 \\
D3 & 价格版本与模型价格 & pricing\_versions, model\_price\_rules, pricing\_rerate\_scopes & 新建版本而非修改 \\
D4 & 用户·密钥·归属期 & telegram\_users, api\_keys, key\_ownership\_periods 等 & 按密钥与时间点查询归属 \\
D5 & 账期·成本快照·账单 & billing\_cycles, cycle\_*, statements, statement\_lines 等 & 关账后只读 \\
D6 & 审计日志 & audit\_logs, reconciliation\_runs & 只追加 \\
\bottomrule
\end{longtable}

\subsubsection{处理过程}

\begin{longtable}{@{}p{0.14\textwidth}p{0.2\textwidth}p{0.58\textwidth}@{}}
\caption{处理过程说明}\\
\toprule 加工 & 输入/输出 & 处理逻辑 \\ \midrule\endfirsthead
\toprule 加工 & 输入/输出 & 处理逻辑 \\ \midrule\endhead
P1 用量同步 & 增量调用记录 / D1 & 校验源表结构指纹；读取检查点之后的记录；解析失败写死信；以“插入且忽略冲突”去重写入；推进检查点 \\
P2 调用计价 & D1、D3、D4 / D2 & 匹配价格模型；确定服务档位与上下文分段；计算等效费用；按资源池规则与归属期确定资源池和用户；幂等写入 \\
P3 密钥管理 & 操作请求 / D4、D6 & 先在网关增删密钥，再写数据库，失败时回滚网关；关闭旧归属期并开启新归属期；吊销使用该密钥的会话 \\
P4 价格管理 & 价格数据 / D3、D6 & 生成新版本并退役旧版本；未关账账期切换到新版本；未变化模型的计价结果直接复制，变化模型重新计价 \\
P5 账期结算 & D2、D3、上游账号 / D5 & 计算各分组成本，扣除按量密钥收入，梯度加权后用最大余数法分摊到分，生成账单与明细；关账时定稿 \\
P6 统计查询 & D1、D2、D5 / 报表 & 按角色限定数据范围，在数据库中完成聚合、分页与分位数计算 \\
\bottomrule
\end{longtable}

\section{总体设计}

\subsection{体系结构}

系统采用 B/S 与消息机器人并行的三层结构（图~\ref{fig:arch}）。表示层是浏览器中的单页应用和 Telegram 客户端；应用层由三个进程组成——Web 服务、Telegram 机器人、后台任务——它们共享同一个业务服务层 \code{BillingService}，保证同一业务规则只有一份实现；数据层是计费数据库。用量库和网关属于外部系统：前者被只读挂载，后者通过 HTTP 管理接口访问。

\begin{figure}[H]
  \centering
  \includegraphics[width=\textwidth]{architecture.pdf}
  \caption{系统体系结构}\label{fig:arch}
\end{figure}

这种划分带来两点好处。第一，所有写操作都经过业务服务层并在同一事务中写入审计日志，无论入口是网页、机器人还是命令行。第二，数据库访问层只依赖 SQLAlchemy 的方言抽象，同一套代码既可运行在单机部署使用的 SQLite 上，也可运行在 PostgreSQL 上；数据库方言的差异集中在少数几处（见 6.2 节），由持续集成在两种数据库上分别跑完全部测试来保证一致。

\subsection{功能模块}

按照需求分析的 F1–F8，系统划分为七个功能模块（图~\ref{fig:modules}）。其中“用量采集计价”由后台任务驱动，其余模块由用户请求驱动。

\begin{figure}[H]
  \centering
  \includegraphics[width=\textwidth]{modules.pdf}
  \caption{功能模块结构}\label{fig:modules}
\end{figure}

\subsection{计费处理流程}

从一次调用到账单的一行，数据要经过图~\ref{fig:pipeline} 所示的处理链。前三步在调用发生后十几秒内由后台任务完成，结果写入计价记录；后五步在账期预览或关账时执行。

\begin{figure}[H]
  \centering
  \includegraphics[width=\textwidth]{pipeline.pdf}
  \caption{从调用记录到账单的处理流程}\label{fig:pipeline}
\end{figure}

\subsection{概念结构设计}

\subsubsection{实体识别}

从数据流图的数据存储出发，识别出如下实体：用户、API 密钥、调用记录、价格版本、模型价格、资源池、上游账号、上游分组、梯度规则、账期、账单、账单明细、账单调整。此外还有三个“带属性的联系”需要特别处理：

\begin{itemize}
  \item \textbf{归属期}：用户与密钥之间不是简单的一对多。一把密钥在不同时间段可以属于不同用户，联系本身带有“生效时间”“失效时间”两个属性。若只在密钥上记录“当前主人”，密钥转手后历史调用就无法正确归属。因此把它设计为联系，转换为关系时成为独立的表。
  \item \textbf{计价}：调用记录与价格版本之间是多对多联系。同一条调用在不同价格版本下各有一个计价结果，联系的属性是等效费用、计费档位和调用时刻的归属用户。
  \item \textbf{成本快照}：账期与上游账号之间是多对多联系。联系属性是该账号在该账期折算出的成本、汇率、分组名等，它们是关账时的历史事实。
\end{itemize}

\subsubsection{E-R 图}

图~\ref{fig:er} 给出系统的全局 E-R 图（陈氏表示法，矩形为实体，菱形为联系，椭圆为联系的属性；实体自身的属性较多，在第 4 章以关系模式列出）。

\begin{figure}[H]
  \centering
  \includegraphics[width=\textwidth]{er.pdf}
  \caption{全局 E-R 图}\label{fig:er}
\end{figure}

主要联系的语义与基数如下：
\begin{itemize}
  \item 用户—归属期—API 密钥（1:N，按时间）：同一时刻一把密钥至多属于一个用户，一个用户可以同时拥有多把密钥；
  \item API 密钥—使用—调用记录（1:N）：调用记录通过密钥摘要关联密钥，未登记的密钥产生的调用也要保留；
  \item 价格版本—包含—模型价格（1:N）：模型价格是弱实体，依赖价格版本存在；
  \item 调用记录—计价—价格版本（M:N）；计价结果归入一个资源池（N:1）；
  \item 账期—采用—价格版本（N:1）：账期创建时绑定生效版本，价格更新时开放账期跟随切换，已关账账期保持不变；
  \item 账期—出账—账单（1:N）且用户—出账—账单（1:N）：每个账期每位用户至多一张账单；
  \item 账单—组成—账单明细（1:N）；明细归入一个资源池（N:1）；
  \item 上游账号—归入—上游分组（N:1），上游分组—遵循—梯度规则（N:1）；
  \item 账期—调整—账单调整（1:N），调整针对某一用户。
\end{itemize}

\subsubsection{关键设计决策}

\paragraph{历史数据版本化而不覆盖}
价格、归属、成本、规则都会变化，而账单必须能复现，因此一律“追加新版本并引用版本号”：改价生成新价格版本，转手是关闭旧归属期、开启新归属期，关账时把规则与成本复制为快照。

\paragraph{计价与分摊分离}
计价只依赖单次调用和价格，可在调用后立即增量完成；分摊依赖整个账期，只在预览或关账时对计价记录聚合一次。

\paragraph{金额全部用整数}
货币量以 64 位整数存储（纳美元、分、百万分之一），既无浮点误差，也回避了 \code{NUMERIC} 在两种数据库上的差异；纳美元精度下可表示约 92 亿美元。

\section{数据库详细设计}

\subsection{逻辑结构设计}

\subsubsection{E-R 图向关系模式的转换}

转换遵循以下规则：每个实体转换为一个关系，实体标识符成为主键；1:N 联系把“1”端主键作为“N”端的外键；带属性的 M:N 联系和带时间属性的 1:N 联系转换为独立关系。为便于阅读，下面用下划线表示主键，用 \# 表示外键，并省略次要属性。

\begin{itemize}
  \item 用户（\underline{用户编号}，用户名，是否管理员，注册时间，……）
  \item API 密钥（\underline{密钥编号}，密钥摘要，登录指纹，脱敏值，状态，\#当前归属用户，计费倍率，……）
  \item 归属期（\underline{归属期编号}，\#密钥编号，\#用户编号，生效时间，失效时间，来源，原因）
  \item 调用记录（\underline{记录编号}，\#数据源，源记录号，事件摘要，发生时间，模型，服务档位，密钥摘要，上游账号，各类 Token，是否失败，延迟，……）
  \item 价格版本（\underline{版本编号}，版本名，状态，来源，生效时间）
  \item 模型价格（\underline{\#版本编号，模型名}，默认四价，Priority 四价，Flex 四价，上下文分段价）
  \item 计价记录（\underline{编号}，\#记录编号，\#版本编号，\#资源池，\#归属用户，发生时间，等效费用，计费档位，计算明细）
  \item 上游分组（\underline{分组编号}，名称，是否默认，\#梯度规则）
  \item 上游账号成本配置（\underline{上游账号}，\#分组编号，订阅方式，订阅起止，每期费用，汇率）
  \item 账期（\underline{账期编号}，账期名，起止时间，状态，\#价格版本，\#梯度规则，梯度规则快照，关账时间，关账人）
  \item 账期上游成本快照（\underline{\#账期编号，上游账号}，账号名，类型，成本金额，汇率，分组快照，订阅快照，……）
  \item 账期分组结算（\underline{\#账期编号，分组编号}，分组名快照，梯度快照，固定成本，按量成本，待分摊成本，……）
  \item 账单（\underline{账单编号}，\#账期编号，\#用户编号，实际用量，加权用量，应付金额，调整金额，是否定稿）
  \item 账单明细（\underline{明细编号}，\#账单编号，\#资源池，分组，实际用量，加权用量，金额，密钥数）
  \item 账单调整（\underline{编号}，\#账期编号，\#用户编号，金额，原因，操作人，时间）
\end{itemize}

完整的逻辑模式共 32 张基本表，按主题分为用户与密钥（6 张）、调用记录（5 张）、价格（5 张）、分摊（3 张）、账单（9 张）、安全与审计（4 张）。核心表之间的主外键关系如图~\ref{fig:schema} 所示，每张表的全部字段见附录 A。

\begin{figure}[H]
  \centering
  \includegraphics[width=\textwidth,height=0.62\textheight,keepaspectratio]{schema.pdf}
  \caption{核心关系模式与参照关系（由模型元数据自动生成）}\label{fig:schema}
\end{figure}

\subsubsection{规范化分析}

以几张代表性的关系为例分析函数依赖。

\paragraph{调用记录 raw\_usage\_events}
候选键为 \code{id}、\code{(source\_id, source\_event\_id)}、\code{(source\_id, event\_hash)}，其余属性完全依赖于每个候选键；“模型 → 价格”不成立（价格还依赖版本），故无部分与传递依赖，属于 BCNF。

\paragraph{模型价格 model\_price\_rules}
各单价依赖于整个主键 \code{(pricing\_version\_id, model)}，属于 BCNF。

\paragraph{归属期 key\_ownership\_periods}
在“同一密钥归属期互不重叠”的约束下，\code{(api\_key\_id, valid\_from\_ms)} 也是候选键；没有非主属性依赖于非键属性，属于 3NF。

\paragraph{账单 statements}
\code{amount\_cents} 在数值上等于明细与调整之和，但关账后它代表“已通知并结清的金额”这一事实，不应随调整记录的误删而改变。保留该属性，并由触发器冻结定稿账单及明细，避免更新异常。

\paragraph{快照属性}
\code{cycle\_groups.group\_name}、\code{billing\_cycles.tiers\_json} 等看似重复主数据，似乎构成传递依赖“(账期, 分组) → 分组 → 分组名”。但分组可以改名或删除、规则可以修改，“分组编号 → 分组名”跨账期并不成立；快照记录的是关账时的取值，依赖于 \code{(cycle\_id, group\_id)} 整体，不构成传递依赖。

\paragraph{计价记录 rated\_events}
\code{occurred\_at\_ms} 可由 \code{raw\_event\_id} 推出，是为“版本 + 时间区间”聚合查询有意保留的冗余（见第 7 章 Q1、Q2）；调用记录不可修改，故不会不一致。\code{telegram\_user\_id} 则是计价时按归属期固化的事实，并非冗余。

综上，全部基本表达到 3NF，多数达到 BCNF；仅有的冗余是历史快照或只读数据上的性能冗余，一致性由约束与触发器保证。

\subsection{物理结构设计}

\subsubsection{数据类型选择}

\begin{itemize}
  \item 所有时间统一为 \code{BIGINT} 毫秒时间戳，不使用 \code{TIMESTAMP}：一是可以直接与用量库对齐，二是避免时区在两种数据库中的语义差异。按时区展示由应用层和视图完成。
  \item 金额与费用使用 \code{BIGINT} 定点整数（见 3.4.4 节），倍率与汇率使用 \code{BIGINT} 百万分之一。
  \item 结构多变、只整体读写的数据（梯度规则、上下文分段价、计算明细、审计前后值）用 \code{TEXT} 存 JSON，不参与查询条件。
  \item 摘要类字段使用 \code{VARCHAR(64)}，名称类字段按业务上限设置长度。
\end{itemize}

\subsubsection{索引设计}

PostgreSQL 中共建立 61 个索引，其中主键 32 个，其余为唯一约束、外键查询和统计查询所需的索引。主要的非主键索引如表~\ref{tab:indexes} 所示。

\begin{longtable}{@{}p{0.3\textwidth}p{0.3\textwidth}p{0.32\textwidth}@{}}
\caption{主要索引设计}\label{tab:indexes}\\
\toprule 索引 & 列 & 服务的查询 \\ \midrule\endfirsthead
\toprule 索引 & 列 & 服务的查询 \\ \midrule\endhead
唯一 (source\_id, source\_event\_id) & raw\_usage\_events & 同步去重、死信重试定位 \\
唯一 (source\_id, event\_hash) & raw\_usage\_events & 同一事件重复导出时的去重 \\
idx\_raw\_events\_time & (occurred\_at\_ms) & 站点状态、时间窗统计 \\
idx\_raw\_events\_key\_time & (api\_key\_hash, occurred\_at\_ms) & 某密钥请求明细，按时间倒序分页 \\
idx\_raw\_events\_missing\_* & (source\_id, source\_event\_id) WHERE 列 IS NULL & 部分索引：只索引待回填的行，回填完成后几乎不占空间 \\
唯一 (raw\_event\_id, pricing\_version\_id) & rated\_events & 计价幂等（冲突即跳过）\\
idx\_rated\_version\_time\_owner\_pool & (version, time, user, pool) & 账期聚合：版本 + 时间区间 + 分组键 \\
idx\_rated\_user\_time & (telegram\_user\_id, occurred\_at\_ms) & 单用户用量查询 \\
idx\_rated\_cycle & (occurred\_at\_ms, pool\_id, user) & 按时间与资源池的聚合 \\
idx\_ownership\_key\_period & (api\_key\_id, valid\_from\_ms, valid\_to\_ms) & 某时刻密钥归属的查找 \\
idx\_api\_keys\_owner\_status & (current\_owner\_id, status) & 用户的有效密钥列表 \\
uq\_pricing\_versions\_single\_active & (status) WHERE status = 'active' & 部分唯一索引：至多一个生效版本 \\
uq\_upstream\_account\_groups\_single\_default & (is\_default) WHERE is\_default & 部分唯一索引：至多一个默认分组 \\
idx\_manual\_usage\_cycle\_group\_pool\_user & (cycle, group, pool, user) & 账期手工用量汇总 \\
\bottomrule
\end{longtable}

索引设计遵循三条原则：（1）复合索引的列顺序为“等值条件列在前、范围条件列在后、分组列随后”，例如 \code{(pricing\_version\_id, occurred\_at\_ms, telegram\_user\_id, pool\_id)}；（2）写入量最大的两张表（调用记录、计价记录）的索引控制在五六个以内，避免拖慢同步；（3）用部分索引表达“只关心少数行”的需求，既省空间又能表达约束。

\subsubsection{视图设计}

为报表与只读角色提供 5 个视图。视图隐藏了密钥摘要、登录指纹等敏感列，并把整数金额换算为元和美元。

\begin{table}[H]
\centering\small
\caption{视图}
\begin{tabularx}{\textwidth}{@{}lX@{}}
\toprule
视图 & 内容 \\ \midrule
v\_user\_cycle\_bills & 每位用户每个账期的账单：实际用量、加权用量、调整、应付金额、是否定稿 \\
v\_statement\_line\_details & 账单明细及其资源池名、分组名 \\
v\_cycle\_group\_summary & 每个账期每个分组的成本核对：总成本、按量收入、待分摊、已分摊、盈余、未分摊 \\
v\_daily\_user\_model\_usage & 生效价格版本下按日（东八区）、用户、模型统计的请求数、失败数、Token 与费用 \\
v\_api\_key\_owners & 密钥与当前归属用户，不含摘要与指纹 \\
\bottomrule
\end{tabularx}
\end{table}

\subsubsection{存储与运行参数}

PostgreSQL 部署时，调用记录和计价记录这两张大表保持默认的 \code{fillfactor = 100}：它们只追加、几乎不更新，页面可以填满。后台任务批量写入后定期依赖 autovacuum 更新统计信息。单机 SQLite 部署时，连接级设置包括：WAL 日志模式（读写并发）、\code{foreign\_keys=on}（启用参照完整性检查）、\code{synchronous=normal}，以及让大排序的临时数据落盘的 \code{temp\_store=file}。

\section{完整性与安全性设计}

\subsection{完整性设计}

\subsubsection{实体完整性}

32 张基本表都定义了主键。按语义分为三类：
\begin{itemize}
  \item 代理键：自增整数，如调用记录、账单；
  \item 自然键：如用户编号直接使用 Telegram 用户 ID，会话主键使用令牌摘要；
  \item 复合键：如模型价格 \code{(pricing\_version\_id, model)}、账期分组结算 \code{(cycle\_id, group\_id)}。
\end{itemize}
另外，业务上需要唯一的属性都声明了 \code{UNIQUE}，包括账期名、价格版本名、密钥摘要、登录指纹，以及“每个账期每位用户一张账单” \code{(cycle\_id, telegram\_user\_id)}、“每条调用在每个版本下只计价一次” \code{(raw\_event\_id, pricing\_version\_id)}。

\subsubsection{参照完整性}

共声明 41 个外键，删除与更新均采用默认的 \code{NO ACTION}（拒绝），不使用级联删除：账务数据被引用就意味着已被账单或计价使用，级联删除会破坏可复现性。需要停用的主数据采用软删除（如规则置 \code{active = false}、密钥置 \code{revoked}）。唯一刻意不建外键的是账期快照中的 \code{group\_id}：分组允许删除，而已关账账期仍要保留当时的分组编号与名称。

\subsubsection{用户定义完整性}

\paragraph{非空与默认值}
必填属性声明 \code{NOT NULL}，计数类属性默认为 0，布尔标志默认为 \code{false}。

\paragraph{CHECK 约束}
共 27 个 CHECK 约束，分为四类：
\begin{itemize}
  \item 状态枚举：密钥状态、账期状态、价格版本状态、密钥操作请求状态；
  \item 非负量：8 类 Token 计数、四项基础单价、等效费用、账单明细与按量收费的金额、订阅费用与汇率；
  \item 区间有序：账期结束晚于开始，归属期失效不早于生效，订阅结束晚于开始；
  \item 非零：手工用量不为零。
\end{itemize}

\paragraph{部分唯一索引}
“同一时刻至多一个生效的价格版本”和“至多一个默认上游分组”无法用普通 UNIQUE 表达，因为非生效版本可以有多个。用带 \code{WHERE} 条件的唯一索引实现，PostgreSQL 与 SQLite 都支持。

\paragraph{触发器}
有些规则涉及多行、多表或状态转换，声明式约束无法表达，用触发器实现（表~\ref{tab:triggers}）。PostgreSQL 中共 16 个触发器、5 个 PL/pgSQL 触发器函数；SQLite 中由于每个触发器只能对应一种操作，相同的规则展开为 33 个触发器。触发器违例时统一抛出 SQLSTATE \code{23000}（integrity\_constraint\_violation），应用层收到的是与违反 UNIQUE 约束同类的 \code{IntegrityError}。

{\small
\begin{longtable}{@{}p{0.28\textwidth}p{0.19\textwidth}p{0.12\textwidth}p{0.35\textwidth}@{}}
\caption{PostgreSQL 触发器清单（BEFORE 行级触发器，由代码自动生成）}\label{tab:triggers}\\
\toprule 触发器 & 表 & 事件 & 条件与动作 \\ \midrule\endfirsthead
\toprule 触发器 & 表 & 事件 & 条件与动作 \\ \midrule\endhead
\input{generated/triggers.tex}
\end{longtable}
}

这些触发器保证的不变量及其业务含义如下：
\begin{enumerate}
  \item \textbf{审计日志只追加}：任何人都不能修改或删除审计记录；
  \item \textbf{已关账账期冻结}：账期行本身，以及它的上游成本快照、分组结算、调整、手工用量都不能再增删改；
  \item \textbf{定稿账单冻结}：\code{final = true} 的账单、按量收费及其明细不可修改。关账操作本身把 \code{final} 从 \code{false} 改为 \code{true}，触发器的条件是 \code{OLD.final}，所以这一步可以执行；
  \item \textbf{已关账账期使用的价格冻结}：被已关账账期引用的价格版本，其模型价格不能修改或删除；
  \item \textbf{调用与计价记录不可篡改}：调用记录中与计费有关的列（Token、模型、时间、密钥、上游账号等）不可修改，也不可删除；用量库后来补充的“推理强度”“服务档位来源”等可选列允许从空值回填。计价记录不可修改；
  \item \textbf{退役价格版本不可复活}；
  \item \textbf{同一密钥的归属期不重叠}：插入新归属期时检查是否与已有区间相交（区间为左闭右开，首尾相接视为不相交）。
\end{enumerate}

第 7 条需要考虑并发：两个事务同时为同一密钥插入归属期时，各自的 \code{EXISTS} 检查都看不到对方未提交的行。因此触发器函数先对密钥行加 \code{FOR UPDATE} 锁，把同一密钥上的插入串行化：

\begin{lstlisting}[language=PLpgSQL]
CREATE OR REPLACE FUNCTION billing_key_ownership_guard() RETURNS trigger LANGUAGE plpgsql AS $$
BEGIN
    PERFORM 1 FROM api_keys WHERE id = NEW.api_key_id FOR UPDATE;
    IF EXISTS (
        SELECT 1 FROM key_ownership_periods p
        WHERE p.api_key_id = NEW.api_key_id
          AND (p.valid_to_ms IS NULL OR p.valid_to_ms > NEW.valid_from_ms)
          AND (NEW.valid_to_ms IS NULL OR NEW.valid_to_ms > p.valid_from_ms)
    ) THEN
        RAISE EXCEPTION 'integrity: key ownership periods overlap'
            USING ERRCODE = 'integrity_constraint_violation';
    END IF;
    RETURN NEW;
END $$;
CREATE TRIGGER trg_key_ownership_no_overlap BEFORE INSERT ON key_ownership_periods
    FOR EACH ROW EXECUTE FUNCTION billing_key_ownership_guard();
\end{lstlisting}

触发器的设计要满足一个前提：不能误拦正常的业务流程。为此先通读了所有写数据库的代码路径，确认关账后不存在对账期子表的写入、审计日志从不更新等事实，然后把触发器装进测试数据库，在 SQLite 和 PostgreSQL 上运行全部 183 个测试用例，没有一个被误拦。设计过程中也发现，最初“模型价格一律不可修改”的规则过严，会妨碍对尚未结算的价格进行修正。后来把它收窄为“已关账账期引用的价格不可修改”，这也正是“历史账单可核对”真正需要的语义。

\subsection{安全性设计}

\subsubsection{应用层安全}

\paragraph{身份认证}
系统不保存任何密码。成员凭 API 密钥登录，数据库中只存两种不可逆值：
\begin{itemize}
  \item 密钥的 SHA-256 摘要，用于与网关的调用记录匹配；
  \item 用服务端密钥计算的 HMAC 指纹，用于登录查找。
\end{itemize}
这样，即使数据库泄露，攻击者也既无法还原密钥，也无法离线枚举登录指纹。管理员令牌通过环境变量配置，比对时使用常数时间比较。

\paragraph{会话管理}
会话令牌为 32 字节随机数，数据库只保存其 HMAC，Cookie 设置 \code{Secure}、\code{HttpOnly} 和 \code{SameSite} 属性。每次请求都会重新校验：登录所用的密钥被吊销或转让后，会话立即失效；管理员令牌轮换后，旧的管理员会话也随之失效。所有写请求都要校验 CSRF 令牌。登录失败按 IP 限流，60 秒内最多 8 次。

\paragraph{分级授权}
每个接口声明所需的角色（表~\ref{tab:roles}）。只读访客的查询在服务层被强制限定为本密钥的数据。管理接口一律要求管理员会话，并把操作者、操作对象、修改前后值写入审计日志。

\paragraph{凭据隔离}
上游账号的 OAuth 令牌只由网关保管。本系统通过网关的代理接口查询额度，请求中只放占位符，从不读取或返回真实凭据；返回给浏览器的上游账号编号也经过脱敏。

\subsubsection{数据库层安全}

PostgreSQL 部署时，用脚本 \code{deployment/postgresql/roles.sql} 建立四个数据库角色，按最小权限原则授权（表~\ref{tab:dbroles}）。

\begin{table}[H]
\centering\small
\caption{数据库角色与权限}\label{tab:dbroles}
\begin{tabularx}{\textwidth}{@{}lX@{}}
\toprule
角色 & 权限 \\ \midrule
billing\_owner & 不可登录；拥有全部表、视图、函数，执行迁移和 DDL 时通过 \code{SET ROLE} 切换为此角色 \\
billing\_app & 应用服务账号：全部表的 SELECT/INSERT/UPDATE/DELETE；收回审计日志、对账记录的 UPDATE/DELETE 权限，收回调用记录、计价记录、归属期的 DELETE 权限；无 DDL 权限 \\
billing\_report & 报表账号：只读 5 个视图、账单与规则类表、用户表的“编号/用户名/注册时间”三列（列级授权）；不能读取密钥表、会话表、审计日志 \\
billing\_auditor & 审计账号：继承 billing\_report，另可读取审计日志、对账记录、死信与同步检查点 \\
\bottomrule
\end{tabularx}
\end{table}

权限与触发器构成两道防线。即使有人删除了触发器，\code{billing\_app} 也无权修改审计日志；而 \code{billing\_app} 自身也没有删除触发器的权限，因为它不是表的属主。下表是在测试库上以各角色实际执行 SQL 的结果：

\begin{table}[H]
\centering\small
\caption{角色权限验证结果}\label{tab:roletest}
\begin{tabularx}{\textwidth}{@{}llX@{}}
\toprule
角色 & 语句 & 结果 \\ \midrule
billing\_app & 应用启动（建表检查、触发器检查） & 成功（对象已存在时不执行 DDL）\\
billing\_app & \code{INSERT INTO audit\_logs …} & 成功 \\
billing\_app & \code{UPDATE audit\_logs …} & 拒绝：permission denied \\
billing\_app & \code{DELETE FROM raw\_usage\_events} & 拒绝：permission denied \\
billing\_app & \code{CREATE TABLE …} / \code{DROP TRIGGER …} & 拒绝：无模式权限 / 非属主 \\
billing\_report & \code{SELECT … FROM v\_user\_cycle\_bills} & 成功 \\
billing\_report & \code{SELECT username FROM telegram\_users} & 成功（授权列）\\
billing\_report & \code{SELECT * FROM telegram\_users} & 拒绝（含未授权列）\\
billing\_report & \code{SELECT cpamp\_hash FROM api\_keys} & 拒绝 \\
billing\_auditor & \code{SELECT … FROM audit\_logs} & 成功 \\
billing\_auditor & \code{SELECT * FROM web\_sessions} & 拒绝 \\
\bottomrule
\end{tabularx}
\end{table}

\section{数据库建立与系统实现}

\subsection{数据库建立}

数据库结构由 SQLAlchemy 模型声明，用 Alembic 迁移脚本管理演进。项目至今共有 17 个迁移版本，从初始模式逐步加入了管理员会话、梯度规则、手工用量、上游分组与订阅、显式上下文分段价等功能。最新的第 17 个迁移加入了本次课程设计完善的完整性对象：部分唯一索引、触发器和视图。新建一个 PostgreSQL 数据库的步骤如下：

\begin{lstlisting}[language=bash]
createdb billing
export BILLING_DATABASE_URL=postgresql+psycopg://billing_owner@localhost/billing
alembic upgrade head                      # 建表、约束、索引、触发器、视图
psql -d billing -v app_password="'...'" -v report_password="'...'" \
     -v auditor_password="'...'" -f deployment/postgresql/roles.sql   # 角色与授权
\end{lstlisting}

已有数据的 SQLite 部署可以用 \code{cpa-billing copy-database <PostgreSQL URL>} 整库迁移。该命令要求目标库为空：先按当前模型建表并写入与源库相同的迁移版本号，再按外键依赖顺序分批复制 32 张表，然后重置自增序列，最后安装触发器。触发器放在数据复制之后安装，原因是如果先装，已关账账期在补录自身历史数据时就会被“冻结”规则拒绝。CHECK 约束和外键则在复制过程中逐行生效，所以历史数据中的任何违规都会被发现。

迁移 0017 面向已有数据的生产库，写成了防御式：建立部分唯一索引前，先检查历史数据是否已有冲突，若有则跳过并告警，不让整个升级失败；触发器只约束之后的写入，不受历史数据影响。

\subsection{多数据库方言适配}

原系统只在 SQLite 上运行。移植到 PostgreSQL 时，用 PostgreSQL 运行全部测试，暴露出以下几类差异，都在数据访问层集中处理（表~\ref{tab:dialect}）。

\begin{longtable}{@{}p{0.25\textwidth}p{0.33\textwidth}p{0.34\textwidth}@{}}
\caption{SQLite 与 PostgreSQL 的差异及处理}\label{tab:dialect}\\
\toprule 差异 & 现象 & 处理 \\ \midrule\endfirsthead
\toprule 差异 & 现象 & 处理 \\ \midrule\endhead
布尔字面量 & 迁移脚本中 \code{SET flag = 1} 在 PostgreSQL 报类型错误 & 改写为 \code{true}（SQLite 3.23 起同样支持）\\
取回新行主键 & \code{cursor.lastrowid} 在 psycopg 中不可用 & 改为 \code{INSERT … RETURNING id} \\
插入时忽略冲突 & \code{INSERT OR IGNORE} 是 SQLite 语法 & 按方言生成 \code{INSERT … ON CONFLICT DO NOTHING} \\
批量插入行数 & psycopg 批量插入的 \code{rowcount} 为 $-1$ & PostgreSQL 下加 \code{RETURNING id} 后计数 \\
标量最大/最小 & SQLite 的多参数 \code{max(a,b)} 在 PostgreSQL 中是聚合函数 & 自定义 \code{greatest/least} 表达式，按方言编译 \\
整数除法 & SQLAlchemy 对整数 \code{/} 生成 \code{CAST(… AS NUMERIC)}，PostgreSQL 再转整数时四舍五入，小时分桶错位 & 改用地板除 \code{//}，两种数据库都生成整数除法 \\
聚合结果类型 & PostgreSQL 的 \code{SUM(bigint)}、\code{AVG} 返回 NUMERIC，到 Python 中是 \code{Decimal} & 注册 NUMERIC 加载器：整数值转 int，非整数转 float，与 SQLite 一致 \\
JSON 取值 & \code{json\_extract} 仅 SQLite 支持 & PostgreSQL 使用 \code{json\_extract\_path\_text} \\
迁移版本表 & Alembic 默认 \code{VARCHAR(32)}，现有版本号长 39 个字符 & 服务端数据库预建 \code{VARCHAR(64)} 的版本表 \\
部分索引 & 只声明了 \code{sqlite\_where} & 同时声明 \code{postgresql\_where} \\
\bottomrule
\end{longtable}

\subsection{计价算法}

设一次调用的输入 Token 为 $I$，输出为 $O$，缓存读取为 $R_0$，缓存写入为 $W$，网关报告的缓存命中数为 $C$。不同上游对缓存的报告口径不一，因此先统一为“缓存读取量” $R$，再求“未命中输入” $U$：
\[
R = \max\bigl(C - R_0 - W,\ 0\bigr) + R_0,\qquad U = \max(I - R - W,\ 0)
\]
单价按以下顺序确定：
\begin{enumerate}
  \item 选模型：从“实际解析的模型名、请求的模型名、记录的模型名”中，取第一个在生效价格版本里有定价的；
  \item 选档位：请求声明 priority 时取 Priority 档，否则依次采用响应档位、记录档位，缺省为 default 档；
  \item 选上下文段：在该档位的上下文分段价中，取阈值严格小于输入长度的最高一段，整次调用的四项单价都用这一段，不跨段拆分。
\end{enumerate}
记四项单价为 $p_{in}, p_{out}, p_{r}, p_{w}$（纳美元/Token），等效费用为
\[
\text{rated\_weight\_nano\_usd} = U\,p_{in} + R\,p_{r} + W\,p_{w} + O\,p_{out}.
\]
计算明细（所用模型、档位、各项 Token 与单价）以 JSON 形式保存在计价记录中，账单上的任何一分钱都可以追溯到具体调用和具体单价。

计价结果写入时以 \code{(raw\_event\_id, pricing\_version\_id)} 为冲突键执行“插入并忽略冲突”，因此后台任务即使重复执行也不会重复计价。价格更新时，对价格未变化的模型，直接用一条 \code{INSERT … SELECT} 把旧版本的计价结果复制到新版本；只有价格变化的模型，才交给后台任务重新计价。

\subsection{分摊算法}

\paragraph{分组成本}
每个上游分组在账期内的总成本 $T$ 由两部分组成：
\begin{itemize}
  \item 包月账号：订阅费按“订阅周期与账期的重叠时长 / 订阅周期时长”的比例折算；
  \item 按量渠道：该渠道上的等效美元用量乘以人民币/美元汇率。
\end{itemize}
未绑定成员的按量密钥按“用量 × 倍率”单独收费，记为 $M$。成员需要分摊的部分是 $\max(T-M, 0)$。

\paragraph{梯度加权}
梯度规则把实际用量（美元）划分为连续的若干段，每段一个倍率，按累进方式计算加权用量：
\[
B(a) = \sum_k \operatorname{round}\Bigl(\max\bigl(0, \min(a, r_k) - l_k\bigr)\cdot m_k\Bigr)
\]
其中 $[l_k, r_k)$ 为第 $k$ 段区间，$m_k$ 为该段倍率。默认规则为：
\begin{itemize}
  \item 0–300 美元，倍率 1；
  \item 300–800 美元，倍率 0.9；
  \item 800–1400 美元，倍率 0.8；
  \item 1400–2000 美元，倍率 0.7；
  \item 2000 美元以上，倍率 0.6。
\end{itemize}
效果是用量越大，边际单价越低。重度用户摊薄了固定的订阅成本，理应获得一定折扣；同时，因为是累进计算，用量多的人付出的总额仍然更多。

\paragraph{最大余数法}
用加权用量作为权重，把待分摊金额（单位：分）分到每个人。先按比例向下取整，再把剩余的几分钱按余数从大到小依次补给各人，余数相同时按用户编号排序，从而保证各人金额之和严格等于待分摊总额。

\paragraph{算例}
某分组本期待分摊 1000.00 元，三位成员的实际用量分别为 A = 500、B = 150、C = 2500 美元。按默认规则：
\begin{itemize}
  \item $B_A = 300 + 200\times0.9 = 480$；
  \item $B_B = 150$；
  \item $B_C = 300 + 500\times0.9 + 600\times0.8 + 600\times0.7 + 500\times0.6 = 1950$；
\end{itemize}
权重合计 2580。以分为单位：
\begin{itemize}
  \item A：$\lfloor 100000\times480/2580\rfloor = 18604$，余数 1680；
  \item B：$\lfloor 100000\times150/2580\rfloor = 5813$，余数 2460；
  \item C：$\lfloor 100000\times1950/2580\rfloor = 75581$，余数 1020。
\end{itemize}
三者合计 99998，还差 2 分，按余数从大到小补给 B 和 A。最终 A 付 186.05 元、B 付 58.14 元、C 付 755.81 元，合计恰为 1000.00 元。与直接按实际用量比例分摊相比，C 的单价由约 0.317 元/美元降为 0.302 元/美元。

\subsection{应用程序实现}

\paragraph{Web 接口与界面}
Web 服务共 53 个路由：认证与会话 7 个，管理功能约 25 个，其余为账务查询、个人密钥、请求明细、站点状态与上游账号。统计（包括延迟分位数、按小时分桶）、分页和条数上限都在 SQL 中完成，应用层只做格式化。前端单页应用（Vue 3 + Vuetify + ECharts）共 9 个页面：账务总览、请求明细、站点状态、上游账号、排行、价格与规则、我的密钥、用户详情和管理控制台。

\paragraph{Telegram 机器人与命令行}
机器人提供注册、密钥管理（\code{/mykey}、\code{/resetkey}、\code{/revoke}）、用量与账单查询（\code{/usage}、\code{/billing}、\code{/chart}）以及管理员命令。命令行工具提供 \code{init}、\code{sync}、\code{preview}、\code{close-cycle}、\code{reconcile}、\code{sync-prices}，以及本次新增的 \code{backup} 和 \code{copy-database}。

\paragraph{代码规模}
后端 Python 约 1.0 万行，前端约 5400 行，迁移脚本 17 个，后端测试 183 个。

\section{系统测试与性能优化}

\subsection{测试方法与环境}

测试分为自动化测试和性能测试两部分：
\begin{itemize}
  \item \textbf{自动化测试}：后端使用 pytest，前端使用 Vitest。每个后端用例都在独立的空数据库上运行，网关和用量库的外部调用一律被拦截，并用临时 SQLite 文件模拟用量库。测试夹具读取环境变量 \code{BILLING\_TEST\_DATABASE\_URL}，同一套用例可以在 SQLite 或 PostgreSQL 上运行，每个用例前重建 \code{public} 模式。GitHub Actions 中配置了 SQLite、PostgreSQL 17 两个矩阵任务，每次提交都会运行。
  \item \textbf{性能测试}：在 Apple M 系列笔记本（PostgreSQL 17.11 默认配置）上进行。
\end{itemize}

\subsection{功能测试}

后端共 183 个测试用例，覆盖了第 2 章的全部功能需求。表~\ref{tab:functest} 按模块列出代表性用例。

\begin{longtable}{@{}p{0.15\textwidth}p{0.5\textwidth}p{0.27\textwidth}@{}}
\caption{代表性功能测试用例}\label{tab:functest}\\
\toprule 模块 & 用例 & 预期 \\ \midrule\endfirsthead
\toprule 模块 & 用例 & 预期 \\ \midrule\endhead
认证 & 未注册用户的密钥登录；仅持有未绑定密钥的访客访问管理接口 & 拒绝；访客只看到本密钥数据 \\
认证 & 吊销登录所用密钥后继续请求；轮换管理员令牌 & 原会话立即失效 \\
密钥 & 重置密钥后新旧密钥的归属期 & 旧期关闭、新期开启，首尾相接 \\
密钥 & 转移密钥后新产生的调用 & 归属到新主人，历史调用不变 \\
采集 & 用量库同一事件重复出现；字段无法解析 & 只导入一次；进入死信并可重试 \\
计价 & Priority 档 + 长上下文；Flex 档无定价时回退 & 使用对应档位/分段价；回退默认档 \\
计价 & 价格同步后仅部分模型变价 & 未变价模型复制结果，变价模型重算，已关账账期不受影响 \\
分摊 & 多个上游分组独立结算后求和 & 各组总额守恒，用户合计正确 \\
分摊 & 按量密钥收入超过分组成本 & 成员分摊为 0，盈余单列 \\
分摊 & 订阅跨月 & 按天数比例折算到账期 \\
账单 & 手工用量在梯度之前计入、跨账期移动 & 两个账期都重新计算 \\
账单 & 关账后再次预览 & 返回定稿账单，不重算 \\
统计 & 小时分桶边界、历史归属下的用户曲线 & 跨小时的调用落入正确的桶 \\
统计 & 延迟分位数、缓存命中率、自定义时间范围 & 与逐条计算的结果一致 \\
\bottomrule
\end{longtable}

\paragraph{测试结果}
\begin{itemize}
  \item SQLite：183 个用例全部通过；
  \item PostgreSQL：182 个通过，1 个按设计跳过（SQLite 专用的在线备份测试）；
  \item 前端：40 个用例全部通过。
\end{itemize}
另外发现 1 个前端用例的期望值依赖运行机器的时区，在东八区或 UTC 下可以通过。

\subsection{完整性与安全性测试}

完整性测试（\code{tests/test\_integrity.py}）针对 5.1.3 节的每一条不变量和各类 CHECK 约束构造违例操作，断言数据库拒绝执行；同时验证合法操作仍然可以执行，例如首尾相接的归属交接、可选列的回填、关账后读取定稿账单。测试还检查视图返回的金额是否正确、触发器重复安装是否幂等。数据库角色权限的验证结果见表~\ref{tab:roletest}。

\subsection{升级兼容性测试}

生产环境中已经有数据的 SQLite 库要能平滑升级，这是“保证生产可用”的关键。测试方法如下：
\begin{enumerate}
  \item 导出改动前的主分支代码，运行其全部测试并保留每个用例留下的数据库，得到 149 个处于不同业务状态的库。这些状态包括开放账期、已关账账期、多价格版本、密钥转移、手工用量等。
  \item 用新代码对这些库执行 \code{alembic stamp 0016} 和 \code{alembic upgrade head}，然后执行服务启动时的初始化流程。
  \item 逐个检查 \code{PRAGMA integrity\_check}、\code{PRAGMA foreign\_key\_check}，确认 33 个触发器全部安装、5 个视图都能查询。
\end{enumerate}
结果：148 个库全部通过。剩下的 1 个是某个用例只打开过连接、没有建表的空库，不代表任何真实部署。此外还把一个包含已关账账期的库用 \code{copy-database} 复制到 PostgreSQL，目标库的迁移版本、视图中的账单金额都正确，直接修改定稿账单会被触发器拒绝。

\subsection{性能测试}

\subsubsection{测试数据与方法}

测试脚本 \code{scripts/benchmark\_postgresql.py} 在空库上按生产模式建表，然后用 \code{COPY} 批量装入合成数据：
\begin{itemize}
  \item 规模：120 位用户、380 把密钥（其中 20 把未绑定，每 10 把中有 1 把在中途转手）、跨度 90 天的 100 万条调用记录，以及 150 万条计价记录（新旧两个价格版本）；
  \item 分布：模型与用户的调用量都服从 Zipf 型长尾分布，Token 数服从对数正态分布；
  \item 耗时：装载 33.5 秒，随后执行 \code{VACUUM ANALYZE}。
\end{itemize}
对系统中 6 类有代表性的查询，各执行 5 次 \code{EXPLAIN (ANALYZE, BUFFERS)}，取执行时间的中位数。“无索引”一组在事务内 \code{DROP INDEX} 后测量，然后回滚。PostgreSQL 的 DDL 是事务性的，回滚后索引自动恢复。

\subsubsection{测试结果}

\begin{table}[H]
\centering\small
\caption{有无索引时的查询耗时（100 万调用记录，中位数）}\label{tab:bench}
\begin{tabularx}{\textwidth}{@{}lXrrr@{}}
\toprule
查询 & 业务场景 & 有索引/ms & 无索引/ms & 加速比 \\ \midrule
Q1 & 某用户本期按模型的用量与费用 & 3.91 & 29.21 & 7.5 \\
Q2 & 账期内按用户、资源池汇总计价用量 & 42.18 & 44.86 & 1.1 \\
Q3 & 某密钥最近 50 条调用（分页）& 0.02 & 0.23 & 12.3 \\
Q4 & 最近 24 小时按模型统计请求、失败、平均延迟 & 3.10 & 32.47 & 10.5 \\
Q5 & 当前版本下尚未计价的调用（反连接）& 153.54 & 145.31 & 0.9 \\
Q6 & 某时刻每把密钥的归属用户 & 0.23 & 5.31 & 22.6 \\
\bottomrule
\end{tabularx}
\end{table}

\begin{figure}[H]
\centering
\begin{tikzpicture}
\begin{axis}[
  ybar, bar width=9pt, width=0.92\textwidth, height=6.2cm,
  ymode=log, log origin=infty, ymin=0.01, ymax=400,
  ylabel={执行时间 / ms（对数坐标）},
  symbolic x coords={Q1,Q2,Q3,Q4,Q5,Q6}, xtick=data,
  legend style={at={(0.5,1.02)},anchor=south,legend columns=2,draw=none,font=\small},
  ymajorgrids, grid style={gray!25},
  nodes near coords, every node near coord/.append style={font=\tiny,rotate=90,anchor=west},
  point meta=rawy,
]
\addplot[fill=blue!45,draw=blue!60!black] coordinates {(Q1,3.91) (Q2,42.18) (Q3,0.02) (Q4,3.10) (Q5,153.54) (Q6,0.23)};
\addplot[fill=orange!55,draw=orange!70!black] coordinates {(Q1,29.21) (Q2,44.86) (Q3,0.23) (Q4,32.47) (Q5,145.31) (Q6,5.31)};
\legend{有索引, 无索引}
\end{axis}
\end{tikzpicture}
\caption{索引对查询耗时的影响}\label{fig:bench}
\end{figure}

\subsubsection{结果分析}

\paragraph{索引显著有效的查询（Q1、Q3、Q4、Q6）}
\begin{itemize}
  \item \textbf{Q1}：执行计划由“计价记录全表扫描”变为“\code{idx\_rated\_user\_time} 位图索引扫描 + 主键嵌套循环连接”。单用户一个月的数据不到全表的 1\%，耗时从 29 ms 降到 4 ms。
  \item \textbf{Q4}：最近 24 小时的数据只占约 1\%，时间索引把顺序扫描变成范围扫描，耗时下降一个数量级。
  \item \textbf{Q6}：每把密钥只需在复合索引 \code{(api\_key\_id, valid\_from\_ms, valid\_to\_ms)} 上定位一两行，不必再对归属期表排序扫描，加速 22 倍。
  \item \textbf{Q3}：删除 \code{(api\_key\_hash, occurred\_at\_ms)} 索引后，优化器改用时间索引倒序扫描，并逐行过滤密钥。由于测试中这把密钥属于高频用户，很快就凑满 50 行，所以无索引时也只有 0.23 ms。对于低频密钥，这种扫描会退化为接近全表扫描。复合索引的价值在于让最坏情况与最好情况一样快。
\end{itemize}

\paragraph{索引作用不大的查询（Q2、Q5）}
测试也给出了两个“索引没起作用”的结果，原因如下：
\begin{itemize}
  \item \textbf{Q2}：账期聚合本来就要读取一个账期的全部计价记录，约占全表的三分之一。此时顺序扫描比随机读索引更便宜，优化器在有索引时同样选择了顺序扫描，这个选择是正确的。这类查询的优化方向不是索引，而是减少执行次数。系统的做法是：账期开放时按需预览；关账后直接读取持久化的账单，不再聚合。
  \item \textbf{Q5}：以全表反连接的方式查找“未计价记录”，两种情况都走哈希反连接，大约 150 ms。系统在实际运行中并不这样查询，而是在内存中维护每个价格版本已计价到的最大记录号，并配合重新计价范围表中的 \code{max\_raw\_event\_id}，只扫描记录号大于水位线的新记录，每轮只处理新增的数百行。这个测试说明了把“扫描全表的去重”改为“基于单调递增主键的增量处理”的必要性。
\end{itemize}

\subsubsection{SQL 改写优化}

视图 \code{v\_daily\_user\_model\_usage} 按东八区日期分组，日期是由时间戳计算出来的表达式。在视图上用 \code{usage\_date BETWEEN …} 过滤时，条件作用在表达式上，无法下推为时间范围扫描。把查询改写为“先用 \code{occurred\_at\_ms} 的范围条件过滤、再计算日期分组”之后，同一报表的耗时从 7.08 ms 降到 1.11 ms。执行计划也从“计价记录位图扫描 + 价格版本哈希连接 + 调用记录嵌套循环”简化为单表的位图索引扫描。系统的统计接口都按改写后的方式编写，视图只供报表角色做即席查询。

\section{系统运行与维护}

\subsection{部署}

系统以 Docker 镜像发布。推送到主分支后，持续集成先在两种数据库上运行测试，再构建镜像并推送到镜像仓库。服务器上的 Docker Compose 启动三个容器：Web 服务、后台任务和机器人。三者共享同一数据库，用量库以只读方式挂载。升级步骤固定为“备份 → 拉取新镜像 → 用新镜像执行 \code{alembic upgrade head} → 重启服务”，迁移成功后才会切换到新版本的服务。

数据库的选择通过环境变量决定：
\begin{itemize}
  \item 未设置 \code{BILLING\_DATABASE\_URL} 时，使用数据目录下的 SQLite 文件，这是单机部署的默认方式；
  \item 设置为 \code{postgresql+psycopg://…} 时，连接 PostgreSQL。此时服务应以 \code{billing\_app} 角色连接，迁移以属主角色执行。
\end{itemize}
服务每次启动都会检查触发器和视图是否齐全，只补建缺失的对象。PostgreSQL 上用会话级咨询锁（advisory lock）避免三个进程同时执行 DDL；SQLite 上用 \code{IF NOT EXISTS} 保证并发启动的安全。

\subsection{备份与恢复}

\begin{itemize}
  \item \textbf{备份}：执行 \code{cpa-billing backup <目标文件>}。SQLite 使用在线备份接口，PostgreSQL 调用 \code{pg\_dump --format=custom}。两种方式都能在服务运行时得到一致的快照，可以配合 cron 每日执行，并保留最近 30 份。
  \item \textbf{SQLite 恢复}：停止三个服务，用备份文件替换数据库文件（同时删除旧的 \code{-wal} 和 \code{-shm} 文件），然后重新启动。
  \item \textbf{PostgreSQL 恢复}：在新的空库上执行 \code{pg\_restore --no-owner -d billing <备份文件>}，再执行一次 \code{roles.sql} 恢复授权。
  \item \textbf{升级前回退点}：每次升级前都要备份。迁移脚本提供 \code{downgrade}，例如 0017 的回退会删除触发器、视图和部分唯一索引，但不会删除任何业务数据。
\end{itemize}

\subsection{日常运维与监控}

\begin{itemize}
  \item \textbf{同步状态}：管理控制台显示同步检查点的积压条数、最近一次成功时间和最近一次错误；后台任务每轮都会在日志中记录导入和计价条数。
  \item \textbf{死信处理}：用量库中无法解析的记录进入死信表，后台任务每轮会自动重试；管理控制台可以查看仍未解决的死信。
  \item \textbf{对账}：\code{cpa-billing reconcile} 比较用量库、原始记录和计价记录三者的条数，统计未归属、未分配资源池和未计价的事件，并把结果写入对账记录表。账期关账前要求未计价事件为 0，否则必须填写数据质量豁免说明。
  \item \textbf{审计}：所有管理操作都记录操作者、对象以及修改前后的值，审计角色可以独立查询。
  \item \textbf{数据增长}：调用记录每月约增长数十万行。当单表达到数千万行时，可以按发生时间对调用记录和计价记录做范围分区，并把已关账账期的分区设为只读。现有索引都以时间为主要或次要列，可以平滑过渡到分区表。
\end{itemize}

\subsection{总结}

本课程设计在一个真实使用中的共享计费系统上，按软件工程流程完成了从问题定义、需求分析到数据库设计、实现、测试和维护的全过程。数据库设计方面的主要工作和体会如下：

\begin{enumerate}
  \item \textbf{用版本化和快照保证账单可复现。}价格、归属、规则都会变化，这是需求最本质的复杂性；价格版本、归属期和账期快照的建模决定了“每期账单可查询、可核对”能否成立。
  \item \textbf{把不变量交给数据库强制执行。}约束、触发器和角色权限在数据库层再落实一遍应用层规则，才能防住手工 SQL 和未来代码的疏漏；设计时先通读写路径，再用全部测试回归，避免误伤正常流程。
  \item \textbf{3NF 是起点而不是教条。}保留的快照属性是历史事实，冗余列有明确的性能收益，二者的一致性由触发器保证。
  \item \textbf{性能结论以测量为准。}百万级测试既验证了索引，也表明账期全量聚合等场景应靠减少计算次数和增量处理来优化。
  \item \textbf{可移植性需要测试保证。}移植暴露的 10 类方言差异中，整数除法的舍入差异不会报错、只会让统计错位，只有双库运行同一套测试才能发现。
\end{enumerate}

后续可以改进的方向包括：按时间对大表分区；为成员提供账单异议与确认流程；把对账任务纳入定时调度并在异常时自动告警。


\appendix
\ctexset{section={name={附录,},number=\Alph{section}}}
\section{数据库表结构（数据字典）}\label{app:dictionary}

本附录由 \code{tools/gen\_schema\_tex.py} 读取系统的 SQLAlchemy 模型自动生成，类型按 PostgreSQL 方言给出。“空”列中的 √ 表示允许为空；在“键”列中：
\begin{itemize}
  \item PK 为主键，FK 为外键，后接被引用的表；
  \item UQ 为单列唯一约束，UQ* 表示该列参与复合唯一约束。
\end{itemize}

{\small
\setlength{\LTpre}{0.3em}\setlength{\LTpost}{0.6em}
\input{generated/data-dictionary.tex}
}

\section{源代码与交付物说明}

\begin{table}[H]
\centering\small
\begin{tabularx}{\textwidth}{@{}lX@{}}
\toprule
路径 & 内容 \\ \midrule
\code{cpa\_billing/models.py} & 数据库模型：表、约束、索引 \\
\code{cpa\_billing/integrity.py} & 触发器与视图的定义和安装（SQLite / PostgreSQL）\\
\code{cpa\_billing/services.py} & 业务服务层：同步、计价、分摊、账单、统计 \\
\code{cpa\_billing/domain.py} & 梯度加权、最大余数法、订阅折算等纯算法 \\
\code{cpa\_billing/web.py}, \code{bot.py}, \code{cli.py} & Web 接口、Telegram 机器人、命令行 \\
\code{cpa\_billing/maintenance.py} & 备份与跨库复制 \\
\code{migrations/versions/} & 17 个 Alembic 迁移脚本 \\
\code{deployment/postgresql/roles.sql} & 数据库角色与授权 \\
\code{scripts/benchmark\_postgresql.py} & 性能测试脚本 \\
\code{frontend/src/} & Vue 3 前端 \\
\code{tests/} & 后端自动化测试（含完整性与维护测试）\\
\code{docs/course-design/} & 本报告的 \LaTeX{} 源文件、图源文件与生成工具 \\
\bottomrule
\end{tabularx}
\end{table}

\section*{参考文献}
\addcontentsline{toc}{section}{参考文献}
\begin{enumerate}[label={[\arabic*]},leftmargin=2.5em]
  \item 王珊，杜小勇，陈红．数据库系统概论（第 6 版）[M]．北京：高等教育出版社，2023．
  \item Silberschatz A, Korth H F, Sudarshan S. Database System Concepts (7th ed.)[M]. McGraw-Hill, 2019.
  \item PostgreSQL Global Development Group. PostgreSQL 17 Documentation: Triggers, Constraints, Privileges, Indexes[EB/OL]. https://www.postgresql.org/docs/17/
  \item SQLAlchemy 2.1 Documentation[EB/OL]. https://docs.sqlalchemy.org/
  \item Alembic Documentation[EB/OL]. https://alembic.sqlalchemy.org/
  \item SQLite Documentation: CREATE TRIGGER, Partial Indexes[EB/OL]. https://www.sqlite.org/docs.html
  \item Balinski M L, Young H P. Fair Representation: Meeting the Ideal of One Man, One Vote (2nd ed.)[M]. Brookings Institution Press, 2001.（最大余数法）
\end{enumerate}


\end{document}
